% GENERATED FILE — DO NOT EDIT; authoritative prose is Markdown.
% source: companion/orthability-and-the-ground-of-intelligibility.md
% source-sha256: ac1f040cd326877d3a79f7161e6afbe7ef5daf778914665663e15cad66c55a8c
% source-qualification: research-stage-draft
% source-qualification: not-peer-reviewed
% source-qualification: not-externally-peer-reviewed
% source-qualification: not-empirically-validated
% source-qualification: candidate-terminology-not-adopted
% source-qualification: philosophical-conclusions-conditional-on-stated-premises
% source-qualification: engineering-evidence-does-not-support-metaphysical-claims
\documentclass[10pt,letterpaper,twocolumn]{article}
\usepackage{amsmath}
\usepackage{amssymb}
\usepackage{booktabs}
\usepackage{geometry}
\usepackage{hyperref}
\usepackage{microtype}
\usepackage{natbib}
\usepackage{xcolor}
\geometry{letterpaper,margin=0.75in}
\hypersetup{colorlinks=true,linkcolor=blue,urlcolor=blue,citecolor=blue}
\setlength{\emergencystretch}{3em}
\title{Orthability and the Ground of Intelligibility: A Transcendental Argument from Rule-Governed Resolution}
\author{}
\date{}
\begin{document}
\twocolumn[
\maketitle
\textbf{Status: complete philosophical DRAFT (R6, 2026-07-20; fresh-session repository review completed; not external human peer review; not empirically validated). Cross-framework dialectical analysis; not peer reviewed; no empirical validation; nothing here is established by the operational manuscript, and nothing here validates it.} “Cross-framework” names accessibility through stated premises and exact rival routing; it is not a criterion-free, autonomous, or coequal tribunal of truth or soundness. Objective orthability makes every proposed criterion objectively assessable. Claim types are labeled inline as \emph{analytic}, \emph{transcendental}, \emph{explanatory}, \emph{teleological}, \emph{metaphysical}, \emph{revelational}, or \emph{school-internal}; no revelational or school-internal claim occurs in this paper (the creed-internal development is the separate Atharī companion). Supersedes nothing: the earlier claim-ladder assessment and outline are retained as supporting history (\texttt{orthemic-modal-metaphysical-assessment.md}, \texttt{orthemic-companion-paper-outline.md}); Thesis C's disposition is Decision 0008. Objection/reply detail: \href{\detokenize{OBJECTIONS-AND-REPLIES.md}}{\texttt{OBJECTIONS-AND-REPLIES.md}}. Sources: \href{\detokenize{sourcing/R3-COMPANION-SOURCING-LEDGER.md}}{\texttt{sourcing/\allowbreak{}R3-\allowbreak{}COMPANION-\allowbreak{}SOURCING-\allowbreak{}LEDGER.\allowbreak{}md}} (current overlay; start at \href{\detokenize{../docs/sourcing/CURRENT-SOURCING-LEDGER.md}}{\texttt{docs/\allowbreak{}sourcing/\allowbreak{}CURRENT-\allowbreak{}SOURCING-\allowbreak{}LEDGER.\allowbreak{}md}}).


\section{Abstract}
Created systems — perceivers, diagnosticians, parsers, validators — resolve concrete cases into determinate placements under governing standards, and their episodes are genuinely sound or unsound. This paper asks what that plain fact presupposes. It develops (A) a three-axis analysis separating \emph{object level} (state-type, represented standard, case-bound token, execution, placement) from \emph{mode of instantiation} (normative specification vs concrete instantiation) and \emph{normative status} (each status assigned to its proper bearer), with strict soundness a derived, factive status stronger than pathway adequacy; (B) an account of orthing \emph{faculties}, first-order and higher-order, embodied and apprehended, engaging reliabilism and proper functionalism; (C) a transcendental argument from \textbf{orthability} — the antecedent determinability, differentiability, and fitness-for-lawful-resolution of what is — to an underived ground of intelligibility, carefully distinguished from modal-ontological and contingency arguments; (D) a reductio against the wholesale creation or self-assembly of orthability, with its strongest objection stated and adjudicated rather than dismissed; (E) a replacement of the earlier token-essential-modality thesis by an evaluability-precondition thesis; and (F) a conditional argument that the ground, if intelligent, is not intrinsically incapable of determinate disclosure. Rival accounts — modal primitivism, Platonism, structural realism, nominalism, selected-function naturalism, semantic naturalism, divine conceptualism, infinitism, deflationism, and generic classical theism — receive standing treatment, and the paper closes by stating exactly which steps each rival blocks and what remains unresolved. The argument is transcendental and internal to modal-normative discourse; it is not presented as a neutral proof compelling every rational reader.

\vspace{1em}
]

\section{1. The phenomenon, imported without its authority}
The operational manuscript describes an architecture in which concrete occurrences (\emph{orthemmata}) are resolved into repeatable operational state-types (\emph{orthemes}) relative to a declared analysis, by governed processes (\emph{orthing}) whose episodes carry separable verdicts of result correctness and pathway adequacy. From that work this paper imports only the \emph{phenomenon}, never its authority: rule-governed, evaluable resolution happens; some placements are correct and some are not; some procedures are truth-conducive and some merely lucky. (\emph{Analytic/empirical import; firewall:} no engineering result supports any conclusion below — the manuscript's evidence stops at the existence and structure of the practice.)

Vocabulary note — three senses, never interchanged (Decision 0010). This paper distinguishes: \textbf{orthability-L} (local, analysis-relative aptness): $\operatorname{Orthable}(m; A)$ — occurrence $m$ is apt for determinate, correct/incorrect resolution under a declared or possible analysis $A$; \textbf{orthability-O} (objective conditions of evaluability): the identity, difference, modal possibility, truth/error, fittingness, lawful-transition, and success/failure conditions that make any adequate analysis, and any evaluable orthing, possible; \textbf{orthability-R} (created representation/application): a finite practice's observation, evidence, placement, rule, and execution. The earlier formulation "always relative to the resolving practice's standards" described sense L only and is superseded: sense O is not practice-relative, and the transcendental argument concerns O — reached from actual episodes through an explicit bridge (§4.2 step 2), never by promoting the practice-relative predicate to a global object by relabeling. Orthability-O is a family of conditions, not an entity beside concrete reality. The word remains a candidate term of the benchmark-gated vocabulary and does no argumentative work: every occurrence can be replaced by its sense-phrase. Full step map: \href{\detokenize{ARGUMENT-MAP-ORTHABILITY.md}}{\texttt{ARGUMENT-MAP-ORTHABILITY.md}}.


\section{2. Concrete, objective, and sound: three axes, not two noun classes}
A diagnosis exists in the clinician; a parse exists in the parser. In the operational vocabulary these are \textbf{concrete placement tokens}: locally instantiated determinations $\hat p$, real as psychological or computational tokens whatever their accuracy. They are not orthemes: the \emph{ortheme} is the repeatable state-type under which a placement places the case, and what actually obtains is the \textbf{objective profile} — in the operational notation, $O^*(m; A)$, the state-types the case truly instantiates under the declared analysis. A placement is \textbf{correct} ($\operatorname{PlacementCorrect}(\hat p, m, A)$) when it corresponds (or adequately approximates, at the analysis's tolerance) to the objective profile; where this paper uses the stronger word \textbf{sound}, it means correct \emph{and} reached through an adequate pathway (§2.1). So the scheme is not two noun classes ("concrete ortheme" vs "pure ortheme") but a product of independent axes: \emph{object level} (state-type; represented standard; case-bound token; execution; resulting placement) × \emph{mode of instantiation} (normative specification vs concrete instantiation) × \emph{normative status} (each status applied to its proper bearer) × \emph{result status}. A mistaken diagnosis is fully concrete and fully incorrect; its being concrete does not make it correct, and its being incorrect does not make it less concrete — \textbf{concrete and sound are not contraries}. (\emph{Analytic; typing per Decision 0009; full development in \href{\detokenize{CONCRETE-AND-SOUND-REASON.md}}{\texttt{CONCRETE-AND-SOUND-REASON.md}}.})

The same discipline applies one level up. The standard a practice actually uses ("appearance suffices for placement here") is, as a repeatable rule-type, a \textbf{metaortheme} which the practice \emph{represents and applies}; the actor's local grasp of it is a \textbf{represented standard}, adequate or inadequate to the objectively fitting standard for that class of case ($\operatorname{AdequateMetaType}(\mu_{\mathrm{rep}}; A)$); and where the standard is materially bound to this case — this datum plane, this instrument, this calibration, this source set — the bound token is a \textbf{metaorthemma}, not a "concrete metaortheme." A laboratory that validates compounds by color alone has a real, concretely operative represented standard which is inadequately related to the fitting one. Status vocabulary distributes across the levels: the \emph{type} is fitting or misfitting; the \emph{representation} adequate or inadequate; the \emph{token} properly or improperly bound, current or stale; the \emph{execution} faithful or deviant; the \emph{placement} correct or incorrect. A single scalar "sound" flag would hide exactly the independent defects the operational framework was built to keep separate. (\emph{Analytic.})


\subsection{2.1 Sound reasoning is stronger than pathway adequacy}
The operational framework makes pathway adequacy deliberately \textbf{non-factive}: $\operatorname{PathwayAdequate}(e)$ does not entail result correctness, because a sufficiently reliable, properly configured, faithfully executed procedure can produce a rare miss. If "sound reasoning" is used in the strict, truth-directed sense — reasoning that genuinely arrives at truth through adequate means — it is therefore \emph{stronger} than pathway adequacy. This paper fixes the derived, claim-level status:


\begin{quote}
$\operatorname{StrictlySoundReasoning}_q(e) := \operatorname{ReasoningPathAdequate}_q(e) \wedge \operatorname{TokenTruthLinked}_q(e)$

\end{quote}
where $\operatorname{ReasoningPathAdequate}_q(e)$ conjoins only the \textbf{claim-relevant} required verdicts $\operatorname{ReqReason}_q(e) \subseteq \operatorname{ReqPath}(e)$ — those bearing evidentially or causally on claim $q$, derived from its dependencies, evidence, governing tokens, and the analysis's governance rules, never selected after the outcome. (Decision 0011 corrects the earlier definition, which conjoined the whole-episode $\operatorname{PathwayAdequate}(e)$: a claim soundly reached should not become unsound because the operator later routed the case poorly or made a false closure claim, and one episode must be able to carry both sound and unsound claims. $\operatorname{PathwayAdequate}(e)$ remains the whole-episode predicate.) The result is factive (truth-linkage is factive) without collapsing the framework's result/pathway independence. It is a \emph{derived profile over existing verdicts}, not a new primitive verdict. The four cells remain distinct: correct-and-adequate (successful, adequately grounded reasoning); correct-but-defective (lucky or compensating-error result — not strictly sound); incorrect-but-adequate (a justified or procedurally adequate miss — not strictly sound in the factive sense, though \texttt{EX\_ANTE\_JUSTIFIED} may pass); incorrect-and-defective (compound failure). "Sound" therefore cannot simply mean "concrete," "reasonable," or "pathway adequate"; whenever it appears without the strict definition, its bearer and dimensions must be declared. (\emph{Analytic + operational import; Decision 0009.})


\subsection{2.2 Concrete reason and ʿaql ṣarīḥ}
A parallel distinction applies to reason itself, and it has a history. "Concrete reason" versus "ideal reason" is a \emph{modern interpretive contrast} — one the supplied scholarship associates with C. Stephen Evans (\emph{Natural Signs and Knowledge of God}, OUP 2010, verified) and reports Jamie B. Turner applying to Ibn Taymiyyah ("Ibn Taymiyya on theistic signs and knowledge of God," \emph{Religious Studies} 58(3), 2022, verified; the exact phrase pair was not confirmed in the accessible text of either work this pass — an open verification item in the sourcing ledger) — and in any case not any classical author's own fixed vocabulary. Used with that label, "concrete reason" may denote either (a) a historically embodied conception of rationality — a community's operative package of accepted premises, inferential repertoire, semantic conventions, evidential authorities, and habitual exclusions, which is not one state-type but a composite of analysis components, represented standards, policies, and practices, soundly related to the fitting standards in some respects and unsoundly in others — or (b) one actual reasoning episode, an actor's dated, situated reasoning run with its trace, premises, governing tokens, and conclusion. Neither sense is pejorative: concrete reasoning can be sound or unsound; concreteness answers \emph{how it is instantiated}, soundness \emph{whether it succeeds}. In the Taymiyyan reconstruction (El-Tobgui's study of the \emph{Darʾ Taʿāruḍ al-ʿaql wa-l-naql}; reported here as secondary scholarship, with primary loci carried as a verification queue in the sourcing ledger), the load-bearing contrast is instead between \emph{ʿaql ṣarīḥ} — clear, unadulterated reason operating through a sound natural disposition and valid sources — and reason corrupted by inherited doctrine, whim, conjecture, doubt, ulterior motive, habit, or imitation: a normative distinction \emph{within} concretely instantiated reason, not a contrast between the concrete and the sound. That tradition also treats faculties as capable of proper function while fallible, and denies that one universal acquisition route is mandatory for every knower — objective truth does not imply identical pathways, which is precisely the operational framework's separation of the objective profile from route-sufficient apprehension. (\emph{Historical/comparative report, secondary-source-bounded; no school-internal claim is asserted.})

Reconciliation with the operational core (Decision 0001). Analysis-relativity does not soften objectivity: once an analysis is declared, $O^*(m; A)$ is fixed by the occurrence and the analysis, and every actor can be wrong about it. What the multi-axis scheme adds is the \emph{normative} reading of that gap: created minds do not manufacture their epistemic targets; they produce local placements that either conform or fail to conform to an antecedent objective order. Placement correctness is correspondence under the declared analysis — not identity of the local representation with reality in every respect, and not anyone's endorsement; strict soundness (§2.1) additionally demands the adequate, truth-linked pathway. And none of this reifies the types: an ortheme's objectivity resides in the truth of the instantiation fact $\operatorname{Inst}_A(m, o)$, not in the type's subsisting as a separate Platonic entity outside every mind and practice. (\emph{Analytic, given D1.})


\section{3. Faculties of orthing}

\subsection{3.1 First-order and higher-order}
Vision, memory, inference, testimony-assessment, and diagnosis are, in this vocabulary, forms of orthing: each resolves concrete inputs into determinate placements. Abstracting, a \textbf{first-order orthing faculty} is a stable capacity $F_O$ taking orthemmata to (usually correct) concrete placements. Above it operates a \textbf{higher-order (metaorthing) faculty} $F_M$ that resolves questions about the resolving itself: is appearance sufficient here? is this inference form valid? is this witness reliable? is this rule fitted to this class of case? Its outputs are higher-order placements and judgments — about governing types (is this standard adequate?), configuration tokens (is this binding proper and current?), executions (was the rule followed?), and episode adequacy; where recurring judgments warrant a reusable distinction, the faculty may \emph{propose} a new repeatable metaortheme or a revision to the governing repertoire — but a one-off judgment token does not become a repeatable type by being entertained. (\emph{Analytic taxonomy; typing per Decision 0009.})

A further distinction matters. A faculty's governing standards may be \textbf{embodied} — operative in its functioning without being represented (a child recognizes faces under constraints it cannot state) — or \textbf{apprehended} — explicitly grasped, formulated, and evaluated by reflective cognition. Sound method is typically embodied before it is propositionally apprehended, and apprehension is what makes correction, teaching, and audit possible. (\emph{Analytic/empirical.})


\subsection{3.2 Reliabilism, proper function, and what they leave open}
Process reliabilism (Goldman 1979) explains epistemic standing through the truth-conduciveness of belief-forming processes; it faces the generality problem (Conee \& Feldman 1998): which of the many types a process token instantiates fixes the reference class? The operational framework's insistence on \emph{declared} reference classes ($\operatorname{RelSpec}$) is a governance answer, not a metaphysical one — it shows how a practice can discipline the choice, not that nature has made it. (\emph{Explanatory; concession recorded.})

Proper functionalism (Plantinga 1993) adds what bare reliabilism lacks: the notions of proper and improper functioning, a congenial environment, a design plan aimed at truth, objective success conditions, and the possibility of malfunction. This is a \textbf{modern comparative framework}, not Ibn Taymiyya's named theory, and neither the label nor reliable functioning by itself proves a Designer. In the present vocabulary: a faculty is not merely \emph{de facto} truth-tracking; it is \emph{ordered toward} sound resolution — toward correct identification, valid inference, fitting evidential weighting, appropriate classification, correct practical handling, and the recognition and repair of defective method. Orthemology's contribution to this exchange is only to widen "truth-directed" to "sound-resolution-directed": the target is not only true propositions but sound placements and sound standards. (\emph{Explanatory/comparative.})


\subsection{3.3 The teleological step, stated with its price}
It is tempting to argue: reliable, properly functioning, truth-directed faculties point to a prior Knower who knows both the targets and the fitting means of apprehending them. \textbf{This inference is not available from reliability alone.} The naturalized alternative is serious: selected-effect accounts (Millikan 1984) ground function in selection history — a faculty's function is what its ancestors were selected for — and deliver malfunction, normativity of a kind, and truth-directedness of a kind, with no designer. To reach a Designer one must add at least two further premises and defend them:


\begin{itemize}
\item \textbf{T1 (irreducible teleology):} the truth- and soundness-directedness of cognitive faculties is not exhaustively reducible to selection-for-fitness — e.g., because selection explains \emph{usefulness} and underdetermines \emph{aim at objective soundness as such}, or because the normative surplus of "ought to have placed correctly" outruns "was selected because placings of this kind correlated with survival." (\emph{Teleological; contested.})

\item \textbf{T2 (causal adequacy):} whatever grounds a genuine ordering of faculties to objective targets must itself be adequate to both — knowing the targets, knowing the fitting means, and capable of instituting the fit. (\emph{Metaphysical; a strengthened "principle of proportionate causality"; contested.})

\end{itemize}
If T1 and T2 hold, the design conclusion follows with content: the ground would need to know every concrete reality and every true, fitting, or possible relation by which it may be apprehended and governed — in the framework's formal representation: the obtaining profiles and the fitting governing standards under any adequate analysis, without thereby reifying types as independently subsisting objects — and to wisely fit faculties to targets, will their instantiation, and have power to create them: Knowledge, Wisdom, Will, Power, arising here from the proper-function architecture rather than being bolted on. If T1 fails, the selected-function account stands and this route stops at biology. The paper endorses the \emph{conditional} and records both antecedent premises as contested. Plantinga's evolutionary argument against naturalism (2011) is a neighbor of T1, not a substitute for it. (\emph{Teleological/metaphysical, conditional.})


\section{4. The transcendental argument from orthability}

\subsection{4.1 What kind of argument it is not}

\begin{itemize}
\item \textbf{Not a modal ontological argument.} That argument runs from the possible necessary existence of God through a modal system to actuality (◇□G → □G; Oppy 2019 surveys). Here, by contrast, \emph{modality is partly the explanandum}: the argument begins from actual determinacy, truth and error, valid and invalid inference, function and malfunction, fittingness and misfit, possibility and impossibility, successful and failed resolution — and asks what makes those distinctions real and correctly applicable. (\emph{Analytic contrast.})

\item \textbf{Not a standard contingency argument.} The contingency argument \emph{uses} the categories of possibility and necessity on beings; this argument asks why those categories — and their correct application — are available at all (Reichenbach 2022 for the contrast class). (\emph{Analytic contrast.})

\item \textbf{Not an argument from ignorance.} No step has the form "we cannot explain X, therefore G"; every step claims a presupposition relation, which is either exhibited or fails. (\emph{Methodological.})

\item \textbf{Not a psychological argument from comprehensibility.} The premise is not that humans find the world intelligible (a fact about us) but that evaluations of resolution episodes are \emph{true} — that soundness and unsoundness are really instantiated. (\emph{Methodological.})

\end{itemize}
What it is: a \textbf{transcendental argument} in the standard sense (Stern \& Cheng 2023) — from an undisputed actuality to its necessary preconditions — and it inherits the standard Stroud-type limitation (Stroud 1968): such arguments constrain what the \emph{discourse participant} must grant, and a determined sceptic may retreat to denying the starting point or reinterpreting the preconditions as features of our conceptual scheme. This paper's conclusion is accordingly an \emph{internal-coherence} result: given the reality of evaluable resolution, an underived ground is required; it is not a neutral proof against every retreat. (\emph{Transcendental; scope conceded.})


\subsection{4.2 The chain, with each step labeled}

\begin{enumerate}
\item \textbf{Actual orthing.} Evaluable resolution episodes occur. (\emph{Empirical/analytic launch point.})

\item \textbf{From episodes to objective conditions — the L→O bridge, explicit.} Actual episodes are created-practice performances (sense R), evaluable under declared analyses (sense L). The bridge premise, stated rather than smuggled: those evaluations are \emph{true} — some placements really are correct and some really incorrect — and the truth of an evaluation cannot be constituted by the evaluating practice, on pain of losing the correct/incorrect contrast the practice itself trades on. So the possibility of true evaluation presupposes \textbf{orthability-O}: reality antecedently determinable, differentiable, and fit for lawful resolution — else there is nothing to resolve and no difference between resolving and failing. The move runs from the existence of \emph{evaluable} orthing to the objective conditions of evaluability; it never runs from the practice-relative predicate L to a practice-independent global object by relabeling. (\emph{Transcendental, with the bridge premise stated.})

\item \textbf{Determinate normative-modal structure.} Orthability-O includes real distinctions of correctness/error, possibility/impossibility, identity/difference, and fittingness/misfit — the evaluability precondition \textbf{C′} of Decision 0008: every evaluable episode presupposes a determinate space of placements possible-but-not-made and of ways failure was possible. (\emph{Transcendental; replaces the retired token-essential-modality thesis.})

\item \textbf{The self-assembly problem.} This structure cannot coherently be explained as arising, by an intelligible process, from an absolutely unintelligible state (§5 below). (\emph{Transcendental reductio, with its scope limit stated there.})

\item \textbf{An underived ground.} Therefore not all orthability is derived/constructed/created; some locus of it is underived. (\emph{Metaphysical conclusion of 1–4, conditional on their readings.})

\item \textbf{Anti-primitivist and anti-Platonist bridges.} Against brute primitivism: primitivism is not an explanation but the declaration that none exists; it "survives" only at the price of explanatory arrears (§7.1). Against Platonism: a realm of causally inert abstracta grounds neither the \emph{applicability} of the structure to concrete occurrences nor the \emph{fit} of faculties to targets; inert exemplars do not select, institute, or enforce anything (§7.2). These are \emph{explanatory-superiority} claims, not demonstrations. (\emph{Explanatory.})

\item \textbf{The move to an intelligent ground.} What must the underived locus be adequate to? Determinate distinctions, norms, and possibilities are the \emph{kind of thing intellects trade in}; a ground adequate to intelligibility-as-such is most naturally characterized as knowing. This is the argument's second explicitly conditional joint (with T2-style causal adequacy doing the work): if grounding intelligibility requires adequacy to it, and adequacy to meanings and norms is intellectual, the ground is intelligent. A rival may stop at step 5 with an impersonal necessary structure; the paper marks that exit. (\emph{Metaphysical/explanatory; conditional.})

\item \textbf{Contingency and Will.} The actual pattern of instantiation — that these possibilities are realized rather than others — is contingent; a ground of both the necessary structure and its contingent instantiation must select among possibilities: Will. (\emph{Metaphysical; conditional on a selection-requires-selector premise, stated, not smuggled.})

\item \textbf{Actualization and Power.} Selection without efficacy institutes nothing; the ground must be capable of actualizing what it selects: Power. (\emph{Metaphysical; same conditionality.})

\item \textbf{Ordered fittingness and Wisdom.} The common-premise positive inference from objective fittingness or ordered teleology to divine Wisdom is \textbf{not established here}. The cross-framework lane therefore records this joint as held and routes the selected-function, intrinsic-form, value-realist, Euthyphro, and impersonal exits rather than declaring victory. A source-bounded Atharī/Taymiyyan treatment may make a conditional creed-internal inference from its additional premises, but that is a different epistemic lane and its school identity is not itself warrant. (\emph{Teleological joint held in this paper; creed-internal development separated.})

\item \textbf{Communicative orthability.} Among the possibilities reality sustains is \emph{determinate representation and disclosure}: signs can mean, orders of words can carry judgment, standards can be taught. §8 develops what this does and does not show about speech. (\emph{Analytic observation feeding §8.})

\end{enumerate}
The chain's honest summary: steps 1–5 constitute the transcendental core; steps 6–7 are explanatory-superiority arguments; steps 8–10 are classical metaphysical extensions each flagged with its extra premise; step 11 is a bridge to a conditional capacity claim. Nothing later strengthens anything earlier.


\section{5. The self-assembly reductio}

\subsection{5.1 The strongest formulation}
Suppose \textbf{all orthability-O is created} (or assembled, or emergent-from-nothing-intelligible) — the objective conditions of evaluability themselves, not merely some practice's standards (sense L/R facts are of course created with their practices; that is not in dispute). Then orthability-O passes from non-instantiation to instantiation. But for that transition to be \emph{intelligible as a production of orthability}, there must already be: a determinate distinction between orthability and its absence; a real possibility of its being instantiated; conditions under which the actualization succeeds or fails; identity across the proposed transition; and a relation between the cause and the appropriate result. Those are already elements of orthability — determinacy, modal possibility, fittingness, lawful transition, correct actualization. So the creation of \emph{all} orthability presupposes orthability; hence not all orthability can be created. Equivalently: an account saying intelligibility "assembled itself" describes the assembly through an intelligible structure that must already obtain — it has not generated intelligibility from its absence but smuggled intelligibility into the machinery allegedly producing it. And the point is not "orthing occurs, therefore orthability is necessary" — actual things may be contingent; it is that \emph{the coming-to-be of orthability would itself already presuppose orthability}. (\emph{Transcendental reductio.})


\subsection{5.2 The strongest objection}
The opponent has a principled reply with three prongs. (i) \textbf{Deny the model of emergence.} The reductio assumes that any origination of orthability is itself an \emph{orthing-like, rule-governed transition} with success conditions. The naturalist may deny this: fundamental emergence, or a brute initial state, need not be a "production" satisfying intelligibility conditions — asking for the lawful conditions of the coming-to-be of lawfulness may be a category mistake. (ii) \textbf{Accept brute modal structure.} Primitivism (§7.1) simply takes the modal-normative structure as fundamental and uncreated \emph{without} a ground; the reductio's conclusion — "not all orthability is created" — is one the primitivist happily grants while refusing the further move to a grounding \emph{being}. (iii) \textbf{Demote the metalanguage.} That \emph{our account} of the transition uses intelligible notions shows something about explanation, not about the ontology explained; an explanatory metalanguage need not be ontologically prior to its subject matter.


\subsection{5.3 Adjudication}
The reductio is \textbf{decisive against a specific class of opponents and not universally conclusive.} It defeats any account that models the origin of intelligibility as an intelligible, norm-governed construction — including most actual naturalistic construction stories (evolutionary, social-constructivist, conventionalist), which do describe rule-apt transitions and therefore presuppose what they would build. It does not defeat the opponent who bites the bullet of prong (i) or (ii): one may hold that there is no intelligible story of origination to be had (brutalism about the initial state) or that modal structure is primitive. Against those positions the argument's force is comparative, not demonstrative: they purchase consistency by declaring the central fact inexplicable, and §7 argues that is an explanatory cost, not a refutation. Prong (iii) is the weakest: the reductio's premise is not that our \emph{description} is intelligible but that the \emph{transition described} — a passage with identity, possibility, and success conditions — instantiates orthability; an opponent who strips those features has retreated to prong (i). (\emph{Adjudication; the "universally conclusive" claim is expressly not made.})


\section{6. The evaluability thesis (Thesis C replaced)}
The earlier assessment carried Thesis C: "every orthing token is essentially modal," supported by a quus-style rule-individuation argument and facing dispositionalist, deflationist, and minimal-implementation objections. Per Decision 0008 this paper argues from the weaker \textbf{C′}: every \emph{evaluable} token presupposes a determinate possibility space for its evaluation. C′ is all the transcendental chain consumes (step 3), and the objections to C do not touch it: a dispositionalist about token modality still needs the distinction between success and possible failure to be \emph{real} for evaluations to be \emph{true}; a deflationist who keeps evaluation-talk true has granted C′ under another name, and one who abandons its truth has exited the starting point, paying the Stroudian toll of §4.1. The rule-individuation argument is retained in the assessment file as an optional strengthening, explicitly not load-bearing. (\emph{Analytic + decision record.})


\section{7. Rival views}
Full objection/reply development: \href{\detokenize{OBJECTIONS-AND-REPLIES.md}}{\texttt{OBJECTIONS-AND-REPLIES.md}}. Standing summary, with each rival's blocking point:


% breakable-row-table: data-rows=10 total-rendered-characters=2829 max-row-rendered-characters=339
\par\medskip
\hrule height 0.8pt
\smallskip
% breakable-row: 1/10
\noindent\textbf{Rival}: \textbf{Brute modal primitivism}\par
\noindent\textbf{What it grants}: Uncreated modal-normative structure (step 5!)\par
\noindent\textbf{Where it blocks the chain}: Step 7 (no intelligent ground; structure is bedrock)\par
\noindent\textbf{Standing assessment}: Consistent; explanatorily arrears-running — it re-labels the explanandum as primitive. Unrefuted.\par
\smallskip
\hrule height 0.4pt
\smallskip
% breakable-row: 2/10
\noindent\textbf{Rival}: \textbf{Platonism / necessary abstracta}\par
\noindent\textbf{What it grants}: Objective, necessary structure\par
\noindent\textbf{Where it blocks the chain}: Steps 6–7 (ground is inert abstracta)\par
\noindent\textbf{Standing assessment}: Faces the applicability/instantiation problem (why does the concrete \emph{participate}?) and offers no account of faculty-target fit. Unrefuted as ontology; weak as explanation.\par
\smallskip
\hrule height 0.4pt
\smallskip
% breakable-row: 3/10
\noindent\textbf{Rival}: \textbf{Structural realism}\par
\noindent\textbf{What it grants}: Objective structure of the actual world\par
\noindent\textbf{Where it blocks the chain}: Steps 2→3 read modally (structure may be actual-world-only)\par
\noindent\textbf{Standing assessment}: Must still ground the modal force of evaluation (possible-but-not-made placements); tends to collapse into primitivism or Platonism at that point.\par
\smallskip
\hrule height 0.4pt
\smallskip
% breakable-row: 4/10
\noindent\textbf{Rival}: \textbf{Nominalism / conventionalism}\par
\noindent\textbf{What it grants}: Concrete particulars + our classifying practices\par
\noindent\textbf{Where it blocks the chain}: Step 2 (no antecedent determinability; we impose)\par
\noindent\textbf{Standing assessment}: Founders on error: if classification is imposition, unsoundness of an imposition needs a standard the view disallows; sophisticated versions re-import objective similarity, i.e., orthability.\par
\smallskip
\hrule height 0.4pt
\smallskip
% breakable-row: 5/10
\noindent\textbf{Rival}: \textbf{Selected-function naturalism}\par
\noindent\textbf{What it grants}: Real (etiological) function and malfunction\par
\noindent\textbf{Where it blocks the chain}: §3.3 T1 (teleology reduced to selection)\par
\noindent\textbf{Standing assessment}: The serious naturalist option. Answered only conditionally (T1); if T1 fails, the faculty route stops at biology — recorded, not hidden.\par
\smallskip
\hrule height 0.4pt
\smallskip
% breakable-row: 6/10
\noindent\textbf{Rival}: \textbf{Semantic naturalism}\par
\noindent\textbf{What it grants}: Naturalized content (information/teleosemantics)\par
\noindent\textbf{Where it blocks the chain}: Step 11's richer reading; §3.3\par
\noindent\textbf{Standing assessment}: Same shape as selected-function debate, at the level of meaning; unresolved in general philosophy of mind, and this paper does not pretend otherwise.\par
\smallskip
\hrule height 0.4pt
\smallskip
% breakable-row: 7/10
\noindent\textbf{Rival}: \textbf{Divine conceptualism}\par
\noindent\textbf{What it grants}: Intelligent ground of abstracta/logic (Anderson \& Welty 2011)\par
\noindent\textbf{Where it blocks the chain}: Nothing before step 8; a fellow traveler that identifies the structure with divine thoughts\par
\noindent\textbf{Standing assessment}: Nearest theistic neighbor; differences (attributes vs thoughts; the Atharī paper's refusal to make possibility-space an entity) are doctrinal, handled there.\par
\smallskip
\hrule height 0.4pt
\smallskip
% breakable-row: 8/10
\noindent\textbf{Rival}: \textbf{Infinitism}\par
\noindent\textbf{What it grants}: Each level grounded in a prior level, without end\par
\noindent\textbf{Where it blocks the chain}: Step 5 (well-foundedness of constitutive dependence)\par
\noindent\textbf{Standing assessment}: Blocked only by a no-infinite-regress-of-grounding premise; that premise is defended as the norm in grounding theory but is a premise. Recorded.\par
\smallskip
\hrule height 0.4pt
\smallskip
% breakable-row: 9/10
\noindent\textbf{Rival}: \textbf{Deflationism about normativity}\par
\noindent\textbf{What it grants}: The practice of evaluation-talk\par
\noindent\textbf{Where it blocks the chain}: Step 1 read realistically\par
\noindent\textbf{Standing assessment}: If evaluation-talk is kept truth-apt, C′ follows; if not, the position abandons the datum. The retreat is available and priced (§4.1).\par
\smallskip
\hrule height 0.4pt
\smallskip
% breakable-row: 10/10
\noindent\textbf{Rival}: \textbf{Generic classical theism (no further identification)}\par
\noindent\textbf{What it grants}: Steps 1–10\par
\noindent\textbf{Where it blocks the chain}: Nothing; declines step 11's theological development\par
\noindent\textbf{Standing assessment}: Fully compatible terminus of this paper; the identification of the ground with the God of a particular revelation is \emph{revelational/school-internal} and is made only in the Atharī companion.\par
\smallskip
\hrule height 0.8pt
\par\medskip
\textbf{What remains unresolved,} stated plainly: T1 against selected-function accounts; the selector premise at step 8; well-foundedness against infinitism; and the Stroudian retreat at the very start. The chain's conclusion therefore has the status: \emph{the best explanation available to one who takes evaluable resolution realistically}, with the named exits open and priced. (\emph{Assessment.})


\section{8. A ground capable of disclosure}
Suppose the chain through step 7 (intelligent ground), with its conditionality acknowledged. What follows about \emph{speech}? The argument, carefully bounded:


\begin{enumerate}
\item Orthability includes communicative orthability (step 11): determinate representation, fitting signification, ordered expression, reception, and correction are among the possibilities reality sustains — created speakers actually instantiate them. (\emph{Analytic/empirical.})

\item The ground of orthability is adequate to \emph{all} of it — including the possibility-space of meaning and its fitting modes of expression (T2-style adequacy again). (\emph{Metaphysical; conditional.})

\item A knowing, willing, powerful ground adequate to meanings and their fitting expression is not \emph{intrinsically mute}: nothing in its nature excludes intentional determinate disclosure. To deny it the capacity for disclosure while granting it adequacy to communicative orthability would attribute to created speakers a perfection-relevant capacity whose ground lacks even the capacity — inverting the direction of adequacy. (\emph{Metaphysical, with the required perfection/causal-adequacy premise stated: capacities of effects are not surpassing their sustaining ground in kind.})

\end{enumerate}
\textbf{Boundary, stated as strongly as the claim itself:} reason, so far, supports at most a \emph{capacity} for intentional disclosure. Reason alone establishes nothing about whether the ground has actually spoken, in what words, with what letters or sound, on what occasions, or that any particular text is its speech. Every such particular is \emph{revelational} and lies wholly outside this paper. (\emph{Boundary claim.})


\section{9. Conclusion}
From the reality of evaluable, rule-governed resolution, this paper argued transcendentally to an underived ground of intelligibility; explanatorily against primitivist and Platonist restings; conditionally — with every extra premise flagged — to a knowing, willing, powerful, wise ground not intrinsically incapable of disclosure. The two-axis analysis and the faculty architecture stand on their own as philosophy of the operational framework, whatever becomes of the theology. The paper's negative results are as definite as its positive ones: reliability alone yields no designer; the reductio does not defeat brutalism or primitivism outright; the token-essential-modality thesis was not defended but replaced; and nothing here is a neutral proof. Where the argument's conditional joints resolve is, in the end, a matter for continued philosophy — and, beyond §8's boundary, not for philosophy at all.

\emph{A bounded note on a worked case.} A separate, co-developed noetic-diagnostic runtime (\texttt{applications/\allowbreak{}daee-\allowbreak{}epistemics/\allowbreak{}}, inspected read-only) is a system whose entire operation \textbf{presupposes} real distinctions among sound and unsound cognition, proper and improper warrant, fitting and unfitting evidential standards, and truthful and false closure. It is therefore a useful \textbf{worked explanandum} for the proper-function and orthability program of this paper — an actual artifact that runs \emph{as if} such standards obtain. It does \textbf{not} establish their ultimate ground: it presupposes a normative (and, in its own creed-internal layer, theological) analysis and operationalizes it, so it cannot prove the §8 bridge, irreducible teleology, a Necessary Being, or divine attributes, and its convergence with this paper is not independent evidence (shared owner and model lineage). The reason/revelation boundary of §8 stands unchanged. \emph{R7B bounded extension (candidate):} the worked case is deepened into a dynamic account — a candidate concept of \textbf{dynamic orthability} (lawful, evidence-sensitive, correction-capable trajectories under a declared analysis) and a metaphysical ladder whose rungs 4–8 still require the §8 bridge premises — in the companion \texttt{companion/\allowbreak{}dynamic-\allowbreak{}orthing-\allowbreak{}noetic-\allowbreak{}learning-\allowbreak{}and-\allowbreak{}orthability.\allowbreak{}md} and the applications \texttt{applications/\allowbreak{}latent-\allowbreak{}state-\allowbreak{}orthing/\allowbreak{}} (Decision 0024) and \texttt{applications/\allowbreak{}daee-\allowbreak{}epistemics/\allowbreak{}META-\allowbreak{}NOETIC-\allowbreak{}MEMETICS-\allowbreak{}AND-\allowbreak{}DYNAMIC-\allowbreak{}ORTHING.\allowbreak{}md} (Decision 0025). None of it validates the framework or proves any theological rung; the layer map \texttt{docs/\allowbreak{}architecture/\allowbreak{}ORTHEMOLOGY-\allowbreak{}LAYER-\allowbreak{}MAP.\allowbreak{}md} keeps each layer's evidence status explicit.


\bigskip\hrule\bigskip

\section{References}
Sources actually cited in this paper. Bibliographic records live in \href{\detokenize{../references/orthemology.bib}}{\texttt{references/\allowbreak{}orthemology.\allowbreak{}bib}}; per-claim evidence-access statuses remain in the sourcing surfaces (start at \href{\detokenize{../docs/sourcing/CURRENT-SOURCING-LEDGER.md}}{\texttt{docs/\allowbreak{}sourcing/\allowbreak{}CURRENT-\allowbreak{}SOURCING-\allowbreak{}LEDGER.\allowbreak{}md}}) and are not restated here.


\begin{itemize}
\item Anderson, James N., and Greg Welty (2011). "The Lord of Noncontradiction: An Argument for God from Logic." \emph{Philosophia Christi} 13(2).

\item Conee, Earl, and Richard Feldman (1998). "The Generality Problem for Reliabilism." \emph{Philosophical Studies} 89.

\item El-Tobgui, Carl Sharif (2020). \emph{Ibn Taymiyya on Reason and Revelation: A Study of Darʾ taʿāruḍ al-ʿaql wa-l-naql}. Brill.

\item Evans, C. Stephen (2010). \emph{Natural Signs and Knowledge of God: A New Look at Theistic Arguments}. Oxford University Press.

\item Goldman, Alvin I. (1979). "What Is Justified Belief?" In G. Pappas (ed.), \emph{Justification and Knowledge}. Reidel.

\item Millikan, Ruth Garrett (1984). \emph{Language, Thought, and Other Biological Categories}. MIT Press.

\item Plantinga, Alvin (1993). \emph{Warrant and Proper Function}. Oxford University Press.

\item Plantinga, Alvin (2011). \emph{Where the Conflict Really Lies: Science, Religion, and Naturalism}. Oxford University Press.

\item Stern, Robert, and Tony Cheng (2023). "Transcendental Arguments." \emph{Stanford Encyclopedia of Philosophy} (archive-pinned edition per the R3 ledger).

\item Stroud, Barry (1968). "Transcendental Arguments." \emph{Journal of Philosophy} 65(9).

\item Turner, Jamie B. (2022). "Ibn Taymiyya on Theistic Signs and Knowledge of God." \emph{Religious Studies} 58.

\end{itemize}

\twocolumn
\bibliographystyle{plainnat}
\bibliography{../../../references/orthemology}
\end{document}
